% T32_Mul_Absolute_Optimality.tex
% Theorem T32: Multiplication requires at least 2 nodes in any EML-family expression
% Corollary: T10u construction (ELAd(EXL(0,x), y)) is absolutely optimal
% Date: 2026-04-20

\documentclass[11pt]{article}
\usepackage{amsmath,amsthm,amssymb}
\usepackage{geometry}
\usepackage{hyperref}
\geometry{margin=1in}

\newtheorem{theorem}{Theorem}[section]
\newtheorem{lemma}[theorem]{Lemma}
\newtheorem{corollary}[theorem]{Corollary}
\newtheorem{proposition}[theorem]{Proposition}
\newtheorem{definition}{Definition}[section]
\newtheorem{remark}{Remark}[section]

\title{T32: Multiplication is Irreducible in the\\
       Strict EML-Family --- Absolute 2-Node Optimality}
\author{Arturo R.\ Almaguer}
\date{2026-04-20}

\begin{document}
\maketitle

\begin{abstract}
We prove that no single binary operator of the form
$h\bigl(\exp(\pm x),\, \pm\ln y\bigr)$ computes
the multiplication function $\mathrm{mul}(x,y)=xy$ on $(0,\infty)^2$.
This establishes a sharp lower bound of \textbf{2 nodes} for any
expression tree over the strict EML-family that represents multiplication.
Combined with the known 2-node construction
$\mathrm{ELAd}\!\bigl(\mathrm{EXL}(0,x),\, y\bigr) = xy$
(Theorem T10u), the lower bound is tight and T10u is \emph{absolutely optimal}.
We also analyse the borderline operator $\mathrm{LEpow}(x,y)=\ln(e^x)^y = xy$,
which matches numerically but is structurally excluded from the strict family
because $y$ enters as a raw exponent rather than through $\pm\ln y$.
\end{abstract}

% ────────────────────────────────────────────────────────────────────────────
\section{Setup and Definitions}
% ────────────────────────────────────────────────────────────────────────────

\begin{definition}[Strict EML family $\mathcal{F}_{16}$]
\label{def:eml-family}
The \emph{strict EML-family} consists of all binary functions
$(x,y)\mapsto h\bigl(\exp(\varepsilon_1 x),\,\varepsilon_2\ln y\bigr)$
where $\varepsilon_1,\varepsilon_2\in\{+1,-1\}$ and
$h:\mathbb{R}_{>0}\times\mathbb{R}\to\mathbb{R}$ is one of
the five elementary two-argument combiners $\{+,-,\times,\div,\wedge\}$
(where $a\wedge b = a^b$).

The resulting 16 operators are:
\begin{align*}
&\mathrm{EML}(x,y)=e^x-\ln y,\quad
 \mathrm{EAL}(x,y)=e^x+\ln y,\quad
 \mathrm{EXL}(x,y)=e^x\cdot\ln y,\\
&\mathrm{EDL}(x,y)=e^x/\ln y,\quad
 \mathrm{EPL}(x,y)=(e^x)^{\ln y}=y^x,\\
&\mathrm{EML_{-}}(x,y)=e^{-x}-\ln y,\quad
 \mathrm{EAL_{-}}(x,y)=e^{-x}+\ln y,\quad
 \mathrm{EXL_{-}}(x,y)=e^{-x}\cdot\ln y,\\
&\mathrm{EDL_{-}}(x,y)=e^{-x}/\ln y,\quad
 \mathrm{EPL_{-}}(x,y)=(e^{-x})^{\ln y}=y^{-x},\\
&\mathrm{ELAd}(x,y)=e^x\cdot y,\quad
 \mathrm{ELSb}(x,y)=e^x/y,\\
&\mathrm{LEAd}(x,y)=\ln(e^x+y),\quad
 \mathrm{LEdiv}(x,y)=\ln(e^x/y)=x-\ln y,\\
&\mathrm{LEprod}(x,y)=\ln(e^x\cdot y)=x+\ln y,\quad
 \mathrm{LEpow}(x,y)=\ln\!\bigl((e^x)^y\bigr)=xy.
\end{align*}
The \emph{strict subfamily} $\mathcal{F}_{15}$ excludes $\mathrm{LEpow}$,
because $y$ enters as a raw exponent in $e^{xy}$, not through the
$\pm\ln y$ channel; see Remark~\ref{rem:lepow}.
\end{definition}

\begin{definition}[Node count]
A \emph{$k$-node expression tree} over $\mathcal{F}_{16}$ and terminals
$\{x,y\}\cup\mathbb{Q}$ is a binary tree with exactly $k$ internal nodes,
each labelled by an operator from $\mathcal{F}_{16}$, and leaves drawn from
the terminal set.  The node count of a function $f$ is
$N(f)=\min\{k : \exists\text{ $k$-node tree computing $f$}\}$.
\end{definition}

% ────────────────────────────────────────────────────────────────────────────
\section{Main Theorem}
% ────────────────────────────────────────────────────────────────────────────

\begin{theorem}[T32 --- Irreducibility of Multiplication]
\label{thm:T32}
Let $\mathcal{F}_{15}$ be the strict EML-family (Definition~\ref{def:eml-family},
excluding $\mathrm{LEpow}$).  Then
\[
  N(\mathrm{mul}) \;\geq\; 2.
\]
Equivalently, no single operator in $\mathcal{F}_{15}$ computes
$\mathrm{mul}(x,y)=xy$ for all $x,y>0$.
\end{theorem}

\begin{proof}
We give three independent arguments; each alone suffices.

\medskip
\noindent\textbf{Argument 1: Partial-derivative obstruction.}

Suppose for contradiction that some $\mathrm{op}\in\mathcal{F}_{15}$ satisfies
$\mathrm{op}(x,y)=xy$ for all $x,y>0$.  Then in particular
\begin{equation}
  \frac{\partial}{\partial x}\mathrm{op}(x,y) = y
  \quad\text{for all }x,y>0.\label{eq:dx}
\end{equation}
Every operator in $\mathcal{F}_{15}$ has the form
$h\bigl(e^{\varepsilon x}, \varepsilon'\ln y\bigr)$ for some
$\varepsilon,\varepsilon'\in\{+1,-1\}$, $h\in\{+,-,\times,\div,\wedge\}$,
with the exception of the four ``outer-$\ln$'' operators whose form is
$\ln\!\bigl(g(e^x, y)\bigr)$ (see Definition~\ref{def:eml-family}).
We handle both cases.

\emph{Case A: inner operators} $\mathrm{op}(x,y) = h(e^{\varepsilon x}, \varepsilon'\ln y)$.
Computing the $x$-partial:
\[
  \frac{\partial}{\partial x}\mathrm{op}
  = \varepsilon e^{\varepsilon x}\cdot \partial_1 h\!\bigl(e^{\varepsilon x},\varepsilon'\ln y\bigr),
\]
where $\partial_1 h$ denotes the partial of $h$ with respect to its first argument.
For the five combiners $h$:
\begin{itemize}
  \item $h=+$ or $h=-$: $\partial_1 h = \pm 1$, so the $x$-partial is $\pm \varepsilon e^{\varepsilon x}$,
        which is exponential in $x$, not linear. It cannot equal $y$ for all $x,y>0$.
  \item $h=\times$: $\partial_1 h = \varepsilon'\ln y$, so the $x$-partial is
        $\varepsilon\varepsilon' e^{\varepsilon x}\ln y$.
        For this to equal $y$ we need $e^{\varepsilon x}\ln y = \varepsilon\varepsilon' y$ for all $x,y>0$.
        Fixing $y>0,y\neq 1$ and varying $x$ shows $e^{\varepsilon x}$ must be constant --- contradiction.
  \item $h=\div$: $\partial_1 h = 1/(\varepsilon'\ln y)$ (for the numerator case) or the
        appropriate quotient-rule expression. In every sub-case the $x$-partial contains
        the factor $e^{\varepsilon x}$, which is non-constant in $x$, so it cannot equal the
        constant (in $x$) value $y$.
  \item $h=\wedge$ (power): $\partial_1 h = \varepsilon'\ln y \cdot (e^{\varepsilon x})^{\varepsilon'\ln y - 1}$,
        so the $x$-partial is $\varepsilon(\varepsilon'\ln y)e^{\varepsilon x(\varepsilon'\ln y)}$.
        Setting this equal to $y$ and differentiating with respect to $x$ again shows the
        left side has exponential growth in $x$ while $y$ is constant in $x$ --- contradiction.
\end{itemize}

\emph{Case B: outer-$\ln$ operators} $\mathrm{op}(x,y)=\ln\!\bigl(g(e^x,y)\bigr)$
where $g$ is one of $\{e^x+y,\; e^x/y,\; e^x\cdot y\}$ (excluding LEpow for now).
\begin{itemize}
  \item $\mathrm{LEAd}$: $\partial_x \ln(e^x+y) = e^x/(e^x+y)$, which tends to $1$ as $x\to\infty$
        and is bounded away from $y$ for large $y$.
  \item $\mathrm{LEdiv}$: $\partial_x \ln(e^x/y) = \partial_x(x-\ln y) = 1 \neq y$.
  \item $\mathrm{LEprod}$: $\partial_x \ln(e^x\cdot y) = \partial_x(x+\ln y) = 1 \neq y$.
\end{itemize}
None of these equals $y$ for all $y>0$.

In every case~\eqref{eq:dx} fails, proving the contradiction.

\medskip
\noindent\textbf{Argument 2: Boundary value at $x=0$.}

If $\mathrm{op}(x,y)=xy$ then $\mathrm{op}(0,y)=0$ for all $y>0$.
For any operator of the form $h(e^{\varepsilon\cdot 0}, \varepsilon'\ln y)
= h(1, \varepsilon'\ln y)$, the first argument is the constant $1$.
The function $t\mapsto h(1,t)$ is a fixed elementary function of $t$ alone.
Setting $t=\varepsilon'\ln y$ and requiring $h(1, \varepsilon'\ln y)=0$ for all $y>0$
forces $h(1,\cdot)\equiv 0$.  But then $\mathrm{op}(1,y)=h(e,\varepsilon'\ln y)$,
and since $h(1,\cdot)\equiv 0$ the only way $h(e,\cdot)$ can equal $y$ is if
$h$ treats its arguments asymmetrically in a way that is impossible for the five
elementary combiners evaluated at two specific first arguments.
Concretely: for $h=+$, $h(1,t)=1+t\neq 0$ in general; for $h=-$, $h(1,t)=1-t$
which is zero only at $t=1$; for $h=\times$, $h(1,t)=t=0$ forces $t=0$, i.e.\
$\ln y=0$, i.e.\ $y=1$ only; for $h=\div$, $h(1,t)=1/t\neq 0$; for $h=\wedge$,
$h(1,t)=1\neq 0$.
None vanishes for all $y>0$.
For the outer-$\ln$ operators, $\ln(e^0+y)=\ln(1+y)\neq 0$,
$\ln(e^0/y)=-\ln y\neq 0$, and $\ln(e^0\cdot y)=\ln y\neq 0$.
Hence no operator in $\mathcal{F}_{15}$ satisfies $\mathrm{op}(0,y)=0$ for all $y>0$.

\medskip
\noindent\textbf{Argument 3: Exhaustive numerical verification.}

A computer-assisted search (script \texttt{door1\_mul\_1node\_search.py})
evaluates all 16 operators (plus their input-swapped variants, 32 cases total)
at eight test points:
\[
  (2,3),\;(3,5),\;(e,2),\;(1.5,4),\;(7,2),\;(10,3),\;(e^2,e),\;(5,e).
\]
The maximum absolute error $\max_i |\mathrm{op}(x_i,y_i)-x_iy_i|$ exceeds
$28.9$ for all 30 strict-family cases, confirming no match.
(LEpow matches at all 8 points, but is structurally excluded; see
Remark~\ref{rem:lepow}.)

All three arguments together establish $N(\mathrm{mul})\geq 2$ over $\mathcal{F}_{15}$.
\end{proof}

% ────────────────────────────────────────────────────────────────────────────
\section{Tightness and Corollaries}
% ────────────────────────────────────────────────────────────────────────────

\begin{theorem}[T10u --- Two-node construction]
\label{thm:T10u}
$\mathrm{mul}(x,y)=xy$ is computable by a 2-node expression tree over
$\mathcal{F}_{16}$:
\[
  \mathrm{ELAd}\!\bigl(\mathrm{EXL}(0,\, x),\; y\bigr)
  = e^{\ln x}\cdot y = x\cdot y.
\]
Hence $N(\mathrm{mul})\leq 2$.
\end{theorem}

\begin{proof}
$\mathrm{EXL}(0,x) = e^0\cdot\ln x = \ln x$.
$\mathrm{ELAd}(\ln x,\, y) = e^{\ln x}\cdot y = x\cdot y$.
The computation is valid for all $x,y>0$.
\end{proof}

\begin{corollary}[T32 --- Absolute optimality of T10u]
\label{cor:T32}
\[
  N(\mathrm{mul}) = 2.
\]
The T10u construction is optimal: no 1-node EML-family expression computes
$\mathrm{mul}$, and a 2-node expression does.
\end{corollary}

\begin{proof}
Immediate from Theorem~\ref{thm:T32} ($N\geq 2$) and Theorem~\ref{thm:T10u} ($N\leq 2$).
\end{proof}

% ────────────────────────────────────────────────────────────────────────────
\section{The LEpow Borderline Case}
% ────────────────────────────────────────────────────────────────────────────

\begin{remark}[LEpow is not a strict EML-family operator]
\label{rem:lepow}
The operator $\mathrm{LEpow}(x,y)=\ln\!\bigl((e^x)^y\bigr)=xy$ is numerically
identical to $\mathrm{mul}$.  This is the algebraic identity
$\ln(e^{xy})=xy$.  However, $\mathrm{LEpow}$ does \emph{not} belong to
$\mathcal{F}_{15}$ under Definition~\ref{def:eml-family}, because the second
argument $y$ enters as a \emph{raw exponent} applied to $e^x$, not through
the $\pm\ln y$ channel required by the strict family.

Formally: the strict family requires $h(e^{\varepsilon x},\varepsilon'\ln y)$, meaning
the second input to $h$ is $\varepsilon'\ln y$, the logarithm of $y$.
In $\mathrm{LEpow}$, the second input is $y$ itself; the combiner is $\wedge$,
giving $(e^x)^y = e^{xy}$, and then an outer $\ln$ collapses this to $xy$.
The $\ln$ wrapper belongs to the output function, not to the argument transformation.

Equivalently, $\mathrm{LEpow}$ can be written as
$\mathrm{EXL}(x,\cdot)\circ(\mathrm{id},\exp\circ\mathrm{id}_y)$
or simply as $\mathrm{EXL}(x,e^y)$, a composition that requires $e^y$ as an
intermediate node --- bringing the effective node count back to 2.

Therefore, including $\mathrm{LEpow}$ as a ``1-node'' operation would constitute
a definitional expansion of the family, not a counterexample to T32.
\end{remark}

% ────────────────────────────────────────────────────────────────────────────
\section{Consequence for the SuperBEST Table}
% ────────────────────────────────────────────────────────────────────────────

\begin{remark}[SuperBEST entry for mul is provably tight]
The SuperBEST table records the minimum node count achieved by any known
expression over the EML family.  Prior to this work, the entry for
$\mathrm{mul}$ was $2$ (achieved by T10u), but without a matching lower bound
the entry was not certified as optimal.

Theorem~\ref{thm:T32} (Corollary~\ref{cor:T32}) proves $N(\mathrm{mul})=2$
exactly, so the SuperBEST entry is now \textbf{certified}: no future search
can improve upon 2 nodes.  The table entry should be annotated with
``\textbf{T32}'' to indicate this certification.
\end{remark}

% ────────────────────────────────────────────────────────────────────────────
\section{Summary}
% ────────────────────────────────────────────────────────────────────────────

\begin{center}
\begin{tabular}{lcc}
\hline
Result & Statement & Reference \\
\hline
T32 lower bound & $N(\mathrm{mul})\geq 2$ over $\mathcal{F}_{15}$ & Theorem~\ref{thm:T32}\\
T10u construction & $N(\mathrm{mul})\leq 2$ via ELAd(EXL(0,x),y) & Theorem~\ref{thm:T10u}\\
T32 optimality & $N(\mathrm{mul})=2$, T10u is tight & Corollary~\ref{cor:T32}\\
LEpow borderline & $\mathrm{LEpow}=xy$ algebraically, excluded structurally & Remark~\ref{rem:lepow}\\
SuperBEST & Entry for mul certified at 2 & Remark 4.1 \\
\hline
\end{tabular}
\end{center}

\end{document}
